\section{Conclusion}
\label{sec:conclusion}

We studied paragraph-level context selection as a trust-aware scoring
problem, combining six heterogeneous signals covering relevance,
authority, and cross-source consistency. On HotpotQA's distractor
setting, a six-feature logistic regression trained on only one hundred
questions per seed attained Precision@2 of $0.345 \pm 0.032$, a 3.3
absolute-point improvement over the strongest single-signal baseline
SBERT at $0.312 \pm 0.035$, with the sign of the gap consistent on
every seed. Paragraph F1 and supporting-fact F1 improved in step. The
effect we consider most informative is on calibration: the mean score
gap between gold and distractor paragraphs grew from $0.170$ to
$0.349$, more than doubling, so the learned combiner is not just a
better ranker but a better-calibrated confidence source.

The ablation was equally informative. A uniform average of relevance,
log-length, and cross-paragraph consistency, which encodes the
textbook intuition of ``trustworthy equals relevant plus authoritative
plus corroborated,'' collapsed to Precision@2 of $0.080$, below BM25
and far below SBERT alone. The fitted logistic-regression weights
diagnose the cause: paragraph length and cross-paragraph consistency
are learned with negative sign, because distractors in HotpotQA are
length-matched to gold and topically close enough that cross-source
similarity peaks among distractors rather than among the gold pair.
The classical trust heuristics are correct in prior but wrong in sign
for this retrieval regime. The broader conjecture this invites, which we flag for
future work rather than claim here, is that in other distractor-rich
retrieval pipelines the list of signals that carry positive weight
should not be assumed from first principles and would benefit from
per-benchmark calibration against ground truth. The present paper
supplies one concrete existence proof of that failure mode on
HotpotQA; whether the sign-flip generalizes across benchmarks is an
open question.

A few honest caveats. Our evaluation rests on two hundred held-out
questions per seed drawn from the $7{,}405$-question HotpotQA
distractor validation split; the margin over SBERT, while per-seed
consistent, is only slightly larger than a single seed-std, and a
larger-subset or larger-seed study would tighten the confidence
interval. We measured only retrieval-side selection, not downstream
question-answering accuracy with the selected paragraphs. Our trust
signals are lightweight proxies (log-length as a stand-in for source
authority, a regex-based capitalized-token overlap as a stand-in for
question-entity salience) and do not yet consume real provenance
metadata such as publication date, citation count, or source
reputation; such metadata are unavailable for HotpotQA and reserved
for a future open-web benchmark. Finally, the learned classifier
sees only one hundred training examples per seed, so we have not
characterized how calibration scales with additional supervision.

Two directions feel immediate. First, running the same pipeline on
genuinely heterogeneous retrieval pools (open-web RAG with real
publication dates, documented authority proxies, and explicit
source-conflict labels) should reveal whether the sign-flip of the
authority and consistency heuristics is specific to HotpotQA or a
more general property of distractor-rich retrieval. Second, lifting
the trust score from paragraph selection into generation time, for
example as a token-weight gate on the retrieved context, would let us
measure whether better-calibrated selection translates into fewer
hallucinated attributions downstream.
